\section*{Chapter \ref{ch:s2:commodity-spreads} --- Commodity Spreads}

\begin{solution}{exo:s2:commodity-spreads:1}
$(2 \times 2.50 + 3.00) \times 42 / 3 - 80 = 112 - 80 = \$32$ a barrel.
\end{solution}

\begin{solution}{exo:s2:commodity-spreads:2}
$50 - 7 \times 3 = \$29$ per megawatt-hour; at a heat rate of 7.5, $50 - 22.5 = \$27.50$. A less efficient plant earns less on the same prices.
\end{solution}

\begin{solution}{exo:s2:commodity-spreads:3}
$0.022 \times 300 + 0.11 \times 50 - 10 = 6.60 + 5.50 - 10 = \$2.10$ a bushel (oil at 50 cents a pound is \$0.50 a pound, and $0.11$ is 11
pounds in dollars per cent).
\end{solution}

\begin{solution}{exo:s2:commodity-spreads:4}
Product prices follow crude with a lag: retail and wholesale product markets adjust over days or weeks, so a jump in crude first narrows the crack and
a fall first widens it. The real 3-2-1 crack's daily changes have a correlation of $-0.38$ with WTI's; the synthetic crack's $-0.36$ is planted.
\end{solution}

\begin{solution}{exo:s2:commodity-spreads:5}
Brent is assessed at the London close and WTI at New York's, so a day's spread combines prices from different times; a move in oil after London's
close shows up in WTI today and in Brent tomorrow, and the spread's change reverses. The daily changes have an autocorrelation of $-0.36$. A rule
that trades at the day's close captures that reversal, which no one could trade at those prices; entered a day later, its Sharpe ratio falls from 0.85
to 0.32.
\end{solution}

\begin{solution}{exo:s2:commodity-spreads:6}
Crushing capacity: when the crush is wide, processors buy beans and sell meal and oil, crushing more; when it is narrow, they slow down. Their
arbitrage pulls the spread back toward the processing cost, as Simon found for 1985--1995.
\end{solution}

\begin{solution}{exo:s2:commodity-spreads:7}
Sharpe ratios 0.85, 1.05 and 0.65, in the market 73\%, 56\% and 31\% of days. Choosing 1.5 after seeing all three is a small multiple test:
the reported 1.05 is the best of three trials and overstates what the rule would earn next (Book~4's deflated Sharpe ratio applies). The band should
be fixed before the test, or chosen on data before the evaluation period.
\end{solution}

\begin{solution}{exo:s2:commodity-spreads:8}
The spread comes back to transport cost only while transport capacity is spare. In the synthetic market a pipeline bottleneck widened the spread by
\$15 over six months, and the fade lost \$14.1 in those months after earning \$35.3 in four years; doubling the position doubles that loss. The EIA
attributed the record Brent-WTI spread of 2011 to bottlenecks near Cushing. Fading a \omterm{def:s2:commodity-spreads:location}{location spread} needs a view on capacity, and a stop when it
is full.
\end{solution}

\begin{solution}{pb:s2:commodity-spreads:1}
\begin{enumerate}
\item A \omterm{def:s2:commodity-spreads:processing}{processing-spread trade} holds an input against its outputs in the process's ratio, to profit from the margin; a \omterm{def:s2:commodity-spreads:location}{location spread} is one
commodity's price difference between two places; a \omterm{def:s2:commodity-spreads:quality}{quality spread} is the difference between two grades.
\item Processing capacity and its cost (refining, generating, crushing); transport capacity and cost; the cost of upgrading the lower grade.
\item Crude, gasoline and heating oil futures are cointegrated with stationary spreads, and risk arbitrage in three spreads was historically
profitable, though perhaps no longer.
\item Deviations from a seasonal, trending equilibrium were transitory, with reversion to the recent five-day average, and simulated rules would have
been profitable.
\item Mean \$21.30, half-life about 51 trading days, correlation of daily changes with WTI's $-0.38$.
\item In 2011 Brent averaged about \$111 and WTI about \$95, a record spread the EIA attributed to bottlenecks near Cushing; on 20 April 2020 WTI's
spot price was $-\$36.98$ and Brent stood \$54.34 above it.
\item The two prices are assessed at different closing times; changes have an autocorrelation of $-0.36$, and a same-close rule captures a reversal no
one could trade: the fade's Sharpe ratio falls from 0.85 to 0.32 a day later.
\item It rose in 14 of 20 years, by \$6.17 on average, $t$-statistic 2.8, worst year $-\$11.97$.
\item A \$20 mean, a \$5 seasonal cycle peaking in late April, a deviation reverting with a 40-day half-life, and a planted lag behind crude.
\item Crack fade \$15.4 a year, Sharpe ratio 1.05, correlation $-0.12$; season removed 6.2, 0.44, $-0.07$; February--April 7.4, 0.78, $-0.15$;
location 3.8, 0.48, 0.01; quality 2.6, 0.98, 0.01; combination 1.69, $-0.13$.
\item The plain rule fades the season as well as the deviation; without the season, only the deviation is left.
\item It earned \$35.3 over four years, lost \$14.1 in the bottleneck's first six months and \$9.3 over its life, and earned \$8.4 after.
\item The synthetic crack fade earns a Sharpe ratio of 1.05 with a correlation of $-0.12$ to crude's daily moves; on EIA spot data, 0.48 with
$-0.02$ (0.35 without April 2020, 0.24 a day late).
\item Two of its four rules are one bet on the crack's season; the planted effects are strong; and the combination's risk is measured in a sample
with one bottleneck.
\item Reversion that depends on a physical link that can break; legs with their own margin, liquidity and expiry; ratios that are conventions.
\item On futures, not spot, at synchronous closes, with parameters fixed before the test and the season estimated from past years only.
\item When the transport capacity is full: the spread then prices the scarcity of the pipe, not a deviation.
\item The spark spread with a heat-rate view: it needs a model of the plant stack.
\item Both trade many markets on simple signals with equal risk; spreads remove the level of each commodity, which systematic macro keeps.
\item On the physical activity that links the legs: the margin of refining, generating, crushing or moving the commodity.
\end{enumerate}
\end{solution}

\begin{solution}{iq:s2:commodity-spreads:1}
Three barrels of crude against two of gasoline and one of distillate: a rough refinery margin. Refiners hedge it by selling products and buying crude
forward; airlines and product buyers hedge the other side.
\end{solution}

\begin{solution}{iq:s2:commodity-spreads:2}
Check that the legs are recorded at the same time; test the rule entered a day later; use traded futures, not assessments; look at the
autocorrelation of daily changes (a strongly negative one is a sign of asynchronous or noisy prices); and check that the reversion survives costs
and out-of-sample periods.
\end{solution}

\begin{solution}{iq:s2:commodity-spreads:3}
Only after asking why: if pipelines or export capacity are full, the spread prices scarce transport and can keep widening, as in 2011; if it is a
temporary outage or a delivery squeeze, a fade with a stop may pay. The EIA's view is that near-term changes come from pipeline capacity or
production.
\end{solution}

\begin{solution}{iq:s2:commodity-spreads:4}
Limits on each spread's risk and on each leg's outright position and liquidity; stress tests that break each spread's link (bottlenecks, outages,
squeezes) and a common scenario for all of them; limits near expiries and delivery.
\end{solution}

\begin{solution}{iq:s2:commodity-spreads:5}
Store each leg as its own futures series with its own roll calendar; build the spread from the legs at synchronous timestamps; roll the legs
together, or record the roll mismatch; compute costs per leg, and margin per leg or per recognised spread.
\end{solution}

\begin{solution}{iq:s2:commodity-spreads:6}
Split day $d$ into the part before the earlier recording time and the last fraction $f$. The spread is the price's increment over the last $f$ of the
day, $S_d = W_d$. So $\Delta S_d = W_d - W_{d-1}$ and $\Delta S_{d+1} = W_{d+1} - W_d$, with independent increments of variance $f\sigma^2$: the
covariance is $-f\sigma^2$ and the variance $2f\sigma^2$, so the autocorrelation is $-1/2$ whatever $f$. If the true spread also moves, its own changes
add variance without this covariance, and the autocorrelation lies between $-1/2$ and zero, like the real $-0.36$.
\end{solution}
